\ExerciseSection

\begin{enumerate}
\def\labelenumi{\arabic{enumi})}
\item
  设 X 是林德勒夫空间（第 I 章，§ 9，习题 15），H 是 X 上的连续实值函数所成之集，它关于关系 \(\geqslant\) 是有向的，且 H 中函数的下包络 g 连续。证明存在属于 H 的递减函数序列 \((f_{n})\) ，它逐点收敛于 g。
\item
  设 I 是 R 中的紧区间，\((f_{n})\) 是定义在 I 上的单调实值函数的序列，它在 I 中逐点收敛于一个连续函数 g。证明 g 是单调的，且序列 \((f_{n})\) 在 I 中一致收敛于 g。
\end{enumerate}

¶ 3) 设 \(\Gamma\) 是 \(\mathbf{R}\) 的一个单传递同胚群，赋有逐点收敛拓扑。对每个 \(x \in \mathbf{R}\) ，用 \(s_x\) 表示 \(\Gamma\) 中使 \(s_x(0) = x\) 的元素。证明映射 \(x \to s_x\) 是 \(\mathbf{R}\) 到 \(\Gamma\) 上的同胚{[}利用下述事实：若 \(s \in \Gamma\) 使得 \(x < s(x)\) 对某个 \(x \in \mathbf{R}\) 成立，则 \(y < s(y)\) 对一切 \(y \in \mathbf{R}]\) 成立；由此证明 \(\Gamma\) 作为拓扑群同构于 \(\mathbf{R}\) {[}利用 § 3 的习题 17 a) 与第 V 章，§ 3 的定理 1{]}。

\begin{enumerate}
\def\labelenumi{\arabic{enumi})}
\setcounter{enumi}{3}
\item
  设 X 是紧空间，H 是 C(X; R) 的向量子空间，使得 \(|u| \in H\) 对每个 \(u \in H\) 成立。则 X 上的连续实值函数 f 可用属于 H 的函数一致逼近的充分必要条件是：对 X 的每一对点 x, y，f 满足每一条线性关系 \(\alpha f(x) = \beta f(y)\) ，其中 \((\alpha, \beta)\) 满足 \(\alpha\beta \geqslant 0\) 且 \(\alpha g(x) = \beta g(y)\) 对每个函数 \(g \in H\) 成立。{[}应用第 1 号的命题 2，并注意 H 在映射 \(u \to (u(x) u(y))\) 下的像是整个平面 \(R^{2}\) ，或是方程为 \(\alpha X = \beta Y\) （ \(\alpha\beta \geqslant 0\) ）的直线，或者仅由单个点 \((0, 0)\) 组成。{]}
\item
  设 X 是紧空间，H 是 X 上连续实值函数所成之集。设 R 是等价关系`` \(u(x) = u(y)\) （对所有 \(u \in H\)）''。则 X 上的连续实值函数 f 可用 H 中函数的多项式（相应地，无常数项的多项式）一致逼近，当且仅当 f 在模 R 的每个等价类上为常数{[}相应地，当且仅当 f 在模 R 的每个等价类上为常数，且在使 \(u(x) = 0\) （对所有 \(u \in H\) ）成立的一切点 x 处取值为 0{]}。
\end{enumerate}

¶ 6) 设 X 是完全正则空间，用 \(\mathcal{C}^{\infty}(\mathbf{X};\mathbf{R})\) 表示 \(\mathcal{B}(\mathbf{X};\mathbf{R})\) 中由 X 上全体有界连续实值函数组成的赋范子代数。为使 \(\mathcal{C}^{\infty}(\mathbf{X};\mathbf{R})\) 的每个分离 X 的点且包含常值函数的子代数都在 \(\mathcal{C}^{\infty}(\mathbf{X};\mathbf{R})\) 中稠密，必须且只需 X 是紧的。{[}设 \(\beta\mathbf{X}\) 是 X 的斯通–切赫紧化（第 IX 章，§ 1，习题 7）。证明：若 \(\beta\mathbf{X}-\mathbf{X}\) 非空，则存在 \(\mathcal{C}^{\infty}(\mathbf{X};\mathbf{R})\) 的分离 X 的点且包含常值函数的非稠密子代数，为此注意 \(\mathcal{C}^{\infty}(\mathbf{X};\mathbf{R})\) 可等同于 \(\mathcal{C}(\beta\mathbf{X};\mathbf{R}).]\)

\begin{enumerate}
\def\labelenumi{\arabic{enumi})}
\setcounter{enumi}{6}
\tightlist
\item
  设 X 是非紧的完全正则空间。
\end{enumerate}

\begin{enumerate}
\def\labelenumi{\alph{enumi})}
\tightlist
\item
  证明下列性质等价：
\end{enumerate}

\(\alpha)\) \(\mathcal{C}^{\infty}(\mathbf{X};\mathbf{R})\) 中由形如 \(c + f\) 的函数（其中 \(c\) 是常数而 \(f\) 具有紧支集（第 IX 章，§ 4，第 3 号））组成的子代数在 \(\mathcal{C}^{\infty}(\mathbf{X};\mathbf{R})\) 中稠密。

\(\beta)\) 若 \(\beta X\) 是 \(X\) 的斯通–切赫紧化，则 \(\beta X - X\) 由单个点组成。

γ) X 上与它的拓扑相容且使 X 关于它为预紧的一致结构只有一个。

δ) 若 A, B 是 X 的两个完全分离的闭子集（第 IX 章，§ 1，习题 11），则 A、B 之一是紧的。

\(\zeta)\) 若 \(\mathcal{U}\) 是 \(X\) 上与它的拓扑相容的任一一致结构，\(V\) 是 \(\mathcal{U}\) 的任一围域，则存在子集 \(L\) （ \(X\) 的），它的余集是紧的，使得 \(L \times L \subset V\) 。

θ) X 上与它的拓扑相容的一致结构只有一个。{[}注意 α) 蕴含 X 局部紧，因而在 βX 中是开的，由此证明 α) → β)；为证 β) → α)，利用斯通–魏尔斯特拉斯定理。为证 β) → γ)，利用下述事实：βX 的一致结构在 X 上诱导的一致结构，是与 X 的拓扑相容且使 X 关于它为预紧的最细一致结构。为证 γ) → δ)，利用第 IX 章，§ 1，习题 11。为证 δ) → ζ)，首先证明：若 V 是 U 的任一围域，则存在 z ∈ X 使 V(z) 不是紧的；否则 X 关于 U 将是完备的且是仿紧的（第 II 章，§ 4，习题 9），因而是正规的；然后借助 δ) 证明 X 将是可数紧的（第 IX 章，§ 2，习题 14），最后利用第

II 章，§ 4，习题 6。然后把这个结果应用于由 \(f(x, x') \leqslant 1\) 定义的围域 V，其中 \(f\) 是 X 上的连续伪度量，并由 \(\delta\) ) 推出：存在 X 的子集 L，它的余集是紧的，使得 \(L \times L \subset V\) 。最后，为证 \(\zeta) \Longrightarrow \theta\) )，注意 X 上的每个超滤子或者收敛，或者比 X 中相对紧子集的余集所成的滤子更细。{]}

\begin{enumerate}
\def\labelenumi{\alph{enumi})}
\setcounter{enumi}{1}
\item
  证明满足 a) 的等价条件的空间 X 是伪紧的（第 IX 章，§ 1，习题 22）。{[}利用第 IX 章，§ 1 的习题 12，或者注意：若 X 不是伪紧的，则存在 X 的点列 \((x_{n})\) ，它没有聚点，且存在连续映射 \(f: \mathbf{X} \to [\mathrm{o}, \mathrm{i}]\) ，使得 \(f(x_{2n}) = \mathrm{o}\) 且 \(f(x_{2n+1}) = \mathrm{i}\) 。{]}
\item
  设 Y 是非紧的完全正则空间。证明 \(\beta Y\) 中由 \(\beta Y - Y\) 的一个点的余集组成的子空间 X 满足 a) 的条件。
\end{enumerate}

\begin{enumerate}
\def\labelenumi{\arabic{enumi})}
\setcounter{enumi}{7}
\tightlist
\item
  设 \((\mathbf{X}_i)_{1\leqslant i\leqslant n}\) 是紧空间的有限族，\(\mathbf{X} = \prod_{i=1}^{n}\mathbf{X}_i\) 是它们的乘积，且对每个指标 \(i\) ，设 \(\mathbf{F}_i\) 是 \(\mathbf{X}_i\) 的闭子集。证明 \(\mathbf{X}\) 上在每个集 \(\mathbf{F}_i\times \prod_{j\neq i}\mathbf{X}_j\) ( \(1\leqslant i\leqslant n\) )的并上取值为零的连续实值函数，都可用形如
\end{enumerate}

\[
(x _ {1}, \dots , x _ {n}) \rightarrow u _ {1} (x _ {1}) \dots u _ {n} (x _ {n}),
\]

的函数的有限和一致逼近，其中 \(u_{i}\) 是 \(\mathbf{X}_i\) 上在 \(\mathbf{F}_i\) 上取值为零的连续实值函数（ \(1 \leqslant i \leqslant n\) ）。

* 9) 设 \((r_n)_{n \geqslant 1}\) 是区间 \(\mathbf{I} = [\mathbf{o}, \mathbf{i}]\) （属于 \(\mathbf{R}\) ）中互异有理数按某个次序排成的序列。用归纳法定义闭区间序列 \(\mathbf{I}_n \subset \mathbf{I}\) 如下：\(\mathbf{I}_n\) 以点 \(r_{k_n}\) （其下标是不含于 \(\bigcup_{p < n} \mathbf{I}_p\) 的最小下标）为其中点；它的长度 \(\leqslant \mathbf{i}/4^n\) ，且它与任何区间 \(\mathbf{I}_p\) （ \(p < n\) ）都不相交；于是 \((\mathbf{I}_n)\) 是两两不相交的闭区间序列。在乘积空间 \(\mathbf{I} \times \mathbf{R}\) 上如下定义连续实值函数 \(u\) ：对每个整数 \(n \geqslant \mathbf{i}\) ，函数 \(x \to u(x, n)\) 等于 \(\mathbf{i}\) （在 \(\mathbf{I}_n\) 的一个内点处），等于 \(\mathbf{o}\) （在 \(\mathbf{I}_n\) 的外部），且取值于 \([\mathbf{o}, \mathbf{i}]\) 中；另一方面，对每个 \(x \in \mathbf{I}\) ，函数 \(y \to u(x, y)\) 在每个区间 \([n, n + \mathbf{i}]\) 上是仿射线性的。证明 \(u\) 在 \(\mathbf{I} \times \mathbf{R}\) 中不能用形如 \(v(x)w(y)\) 的函数的线性组合一致逼近，其中 \(v\) 在 \(\mathbf{I}\) 上连续，而 \(w\) 在 \(\mathbf{R}\) 上连续且有界。{[}在空间 \(\mathcal{C}^\infty(\mathbf{R}; \mathbf{R})\) （由 \(\mathbf{R}\) 上的连续有界实值函数组成）中，考察部分映射 \(y \to u(x, y)\) （当 \(x\) 跑遍 \(\mathbf{I}\) 时）所成之集；证明在这些函数中存在无限序列 \((u_n)\) ，使得 \(||u_n|| = \mathbf{I}\) 且 \(||u_m - u_n|| = \mathbf{I}\) （只要 \(m \neq n\) ）。由此推出：不存在任何有限维向量子空间 \(\mathbf{E}\) （ \(\mathcal{C}^\infty(\mathbf{R}; \mathbf{R})\) 的），使得每个 \(u_n\) 到 \(\mathbf{E}\) 的距离 \(\leqslant 1/4\) ；否则将存在点的序列 \((v_n)\) （ \(\mathbf{E}\) 中的），使得 \(||v_n|| = 2\) 且 \(||v_n - v_m|| \geqslant 1/2\) （只要 \(m \neq n\) ），这与 \(\mathcal{C}^\infty(\mathbf{R}; \mathbf{R})\) 的每个有限维子空间都是局部紧的这一事实相矛盾。{]} *

\begin{enumerate}
\def\labelenumi{\arabic{enumi})}
\setcounter{enumi}{9}
\tightlist
\item
  利用斯通–魏尔斯特拉斯定理，对紧空间 X 的闭子空间给出乌雷松定理（第 IX 章，§ 4，第 2 号，定理 2）的一个新证明。{[}若 F 是 X 的闭子空间，考察连续映射 \(\mathbf{X} \to \mathbf{R}\) 在 F 上的限制所成之集 H，并注意：若 \(f \in \mathbf{H}\) 且 \(|f(x)| \leqslant a\) 对一切 \(x \in \mathbf{F}\) 成立，则存在函数 \(g \in \mathbf{H}\) ，它在 F 上与 \(f\) 相合，且使得 \(|g(x)| \leqslant a\) 在 X 中成立。{]}
\end{enumerate}

¶ 11) a) 设 X 是完全正则空间，Y 是 X 的闭子空间；映射 \(\varphi: u \to u|\mathbf{Y}\) （由 \(\mathcal{C}^{\infty}(\mathbf{X}; \mathbf{R})\) 到 \(\mathcal{C}^{\infty}(\mathbf{Y}; \mathbf{R})\) 中）是范数 \(\leqslant 1\) 的连续线性映射。若 X 正规，则 \(\varphi\) 是满射的严格态射；又若用 \(\mathbf{X}_{\mathbf{Y}}\) 表示把 Y 的一切点等同起来所得的 X 的商空间（ \(\mathbf{X}_{\mathbf{Y}}\) 由第 IX 章，§ 4 的习题 15 是正规的），则 \(\varphi\) 的核（作为赋范空间）可等同于 \(\mathcal{C}^{\infty}(\mathbf{X}_{\mathbf{Y}}; \mathbf{R})\) 中由在 \(\omega\) （Y 在 \(\mathbf{X}_{\mathbf{Y}}\) 中的典范像）处取值为零的函数所组成的子空间。

\begin{enumerate}
\def\labelenumi{\alph{enumi})}
\setcounter{enumi}{1}
\tightlist
\item
  假设 X 可度量化，且 Y 在 X 中的边界是可数型的。证明 \(\varphi: u \to u|Y\) 是一个具有右逆的严格态射（第 III 章，§ 6，第 2 号，命题 3）。{[}设 \(d\) 是与 X 的拓扑相容的度量；设 \((a_n)\) 是 Fr (Y) 中的点列，在 Fr (Y) 中稠密，并设 \(\mathbf{V}_{n,m}\) 是点 \(x \in \mathbf{X}\) 中使 \(d(x, a_n) < \mathrm{i}/m\) 且 \(d(x, \mathbf{Y}) > \mathrm{i}/2m\) 者所成之集；构造连续映射族 \(f_{n,m}: \mathbf{X} \to [\mathrm{o}, \mathrm{i}]\) ，使 \(f_{n,m}\) 的支集含于 \(\mathbf{V}_{n,m}\) ，且使 \(\sum_{n,m} f_{n,m}(x) = \mathrm{i}\) 对一切 \(x \in \complement Y\) 成立，而族 \((f_{n,m})\) 在 \(\complement Y\) 的每点的某个邻域中一致可和。然后证明：对每个 \(v \in \mathcal{C}^\infty(\mathbf{Y}; \mathbf{R})\) ，函数 \(u\) （它与 \(v\) 在 Y 上相合，且等于 \(\sum_{n,m} v(a_n) f_{n,m}(x)\) （对一切 \(x \in \complement Y\) ））属于 \(\mathcal{C}^\infty(\mathbf{X}; \mathbf{R})\) ，且使得 \(\varphi(u) = v.\) {]} (*)
\end{enumerate}

\begin{enumerate}
\def\labelenumi{\arabic{enumi})}
\setcounter{enumi}{11}
\tightlist
\item
  \begin{enumerate}
  \def\labelenumii{\alph{enumii})}
  \tightlist
  \item
    设 X 是紧空间，\(f: \mathbf{X} \to \mathbf{Y}\) 是 X 到紧空间 Y 上的连续满射。证明映射 \(^{a}f: u \to u \circ f\) （由 \(\mathcal{C}(\mathbf{Y}; \mathbf{R})\) 到 \(\mathcal{C}(\mathbf{X}; \mathbf{R})\) 中）是 \(\mathcal{C}(\mathbf{Y}; \mathbf{R})\) 到 \(\mathcal{C}(\mathbf{X}; \mathbf{R})\) 的一个含单位元的闭子代数上的（保范的）同构。
  \end{enumerate}
\end{enumerate}

\begin{enumerate}
\def\labelenumi{\alph{enumi})}
\setcounter{enumi}{1}
\item
  反之，设 A 是 \(\mathcal{C}(\mathbf{X};\mathbf{R})\) 的含单位元的闭子代数。证明存在连续满射 \(f\) （由 \(\mathbf{X}\) 到一个紧空间 \(\mathbf{Y}\) 上），使得 \(^a f\) 是 \(\mathcal{C}(\mathbf{Y};\mathbf{R})\) 到 A 上的一个同构。{[}若 \((u_{\lambda})_{\lambda \in \mathbf{L}}\) 是在 A 中稠密的一个族，考察连续映射 \(x\to (u_{\lambda}(x))\) （由 \(\mathbf{X}\) 到 \(\mathbf{R}^{\mathrm{L}}\) 中），并利用斯通–魏尔斯特拉斯定理。{]}
\item
  由 b) 推出：对每个序列 \((u_{n})\) （由 \(\mathcal{C}(\mathbf{X};\mathbf{R})\) 的元素组成），存在一个可度量化的紧空间 Y 与一个连续满射 \(f: X \to Y\) ，使得诸 \(u_{n}\) 属于 \(\mathcal{C}(\mathbf{Y};\mathbf{R})\) 在 \(^{a}f\) 下的像。
\end{enumerate}

¶ 13) 设 X 是紧空间，用 A 表示赋范代数 C(X; R)。对 A 的每个子集 M，用 V(M) 表示点 \(x \in X\) 中使 \(u(x) = o\) （对一切 \(u \in M\) ）成立者所成之集。对 X 的每个子集 Y，用 \(\Im(Y)\) 表示元素 \(u \in A\) 中使 \(u(x) = o\) （对一切 \(x \in Y\) ）成立者所成之集。

\begin{enumerate}
\def\labelenumi{\alph{enumi})}
\tightlist
\item
  V(M) 是 X 的闭子空间，而 \(\Im(Y)\) 是 A 的闭理想。若 a 是由 A 的子集 M 生成的 A 的理想，则 \(\mathrm{V}(\mathrm{M})=\mathrm{V}(\mathfrak{a})\) ；V 与 \(\Im\) 是逆序映射（关于包含序关系），并且我们有 \(V\{o\}=X\) ；\(\mathrm{V}(\mathrm{A})=\varnothing\) ；对 A 的每个单位 u 有 \(\mathrm{V}(\{u\})=\varnothing\) ；
\end{enumerate}

\[
\mathrm{V} \left(\bigcup_ {\lambda \in \mathrm{L}} \mathrm{M} _ {\lambda}\right) = \mathrm{V} \left(\sum_ {\lambda \in \mathrm{L}} \mathrm{M} _ {\lambda}\right) = \bigcap_ {\lambda \in \mathrm{L}} \mathrm{V} (\mathrm{M} _ {\lambda})
\]

对 A 的子集的每个族 \((\mathbf{M}_{\lambda})_{\lambda \in \mathbf{L}}\) 成立；\(\mathbf{V}(\mathbf{M}. \mathbf{M}') = \mathbf{V}(\mathbf{M}) \cup \mathbf{V}(\mathbf{M}')\) ；\(\Im(\mathbf{X}) = \{\mathrm{o}\}\) ；\(\Im(\emptyset) = \mathbf{A}\) ；\(\Im\left(\bigcup_{\lambda \in \mathbf{L}} \mathbf{Y}_{\lambda}\right) = \bigcap_{\lambda \in \mathbf{L}} \Im(\mathbf{Y}_{\lambda})\) 对 X 的子集的每个族 \((\mathbf{Y}_{\lambda})_{\lambda \in \mathbf{L}}\) 成立。

\begin{enumerate}
\def\labelenumi{\alph{enumi})}
\setcounter{enumi}{1}
\item
  证明：若 a 是 A 的任一理想，则 \(\Im(\mathrm{V}(a)) = \overline{a}\) ；又若 Y 是 X 的任一子集，则 \(\mathrm{V}(\Im(\mathrm{Y})) = \overline{\mathrm{Y}}\) 。（对第一个关系式，注意 \(\mathrm{V}(\overline{a}) = \mathrm{V}(a)\) ，因而可设 a 是闭的；注意若 x, y 是 \(\mathrm{X} - \mathrm{V}(a)\) 的两个不同的点，则存在 \(u \in a\) 使 \(u(x) \neq u(y)\) ，并利用第 2 号的命题 4。为证第二个关系式，利用乌雷松定理。）
\item
  由 b) 推出：\(x \to \Im(\{x\})\) 是 X 到 A 的极大理想所成之集上的双射，因而这些极大理想都是闭的。A 的一个理想是闭的，当且仅当它是诸极大理想的一个交。环 A 无根。
\item
  证明映射 \(e \to \mathrm{V}(\{e\})\) 是 A 的幂等元所成之集到 X 的既开又闭的子集所成之集上的双射。
\end{enumerate}

一个点 \(x \in \mathbf{X}\) 是孤立点，当且仅当极大理想 \(\Im(\{x\})\) 是主理想。

\begin{enumerate}
\def\labelenumi{\arabic{enumi})}
\setcounter{enumi}{13}
\tightlist
\item
  设 \(\mathbf{X}\) 是拓扑空间。对每个函数 \(u \in \mathcal{C}(\mathbf{X}; \mathbf{R})\) （它使 \(u(x) > 0\) 对一切 \(x \in \mathbf{X}\) 成立），用 \(\mathbf{V}_u\) 表示一切 \(f \in \mathcal{C}(\mathbf{X}; \mathbf{R})\) 中使 \(|f| \leqslant u\) 者所成之集。
\end{enumerate}

\begin{enumerate}
\def\labelenumi{\alph{enumi})}
\item
  证明诸集 \(V_{u}\) 是 o 在 A = C(X; R) 中关于某个与 A 的环结构相容的豪斯多夫拓扑 C 的邻域。
\item
  由 \(\mathcal{C}\) 在 \(\mathcal{C}^{\infty}(\mathbf{X};\mathbf{R})\) 上诱导的拓扑比 \(\mathbf{X}\) 上一致收敛的拓扑细。这两个拓扑相同，当且仅当 \(\mathbf{X}\) 是伪紧的（第 IX 章，§ 1，习题 21）。若 \(\mathbf{X}\) 不是伪紧的，则由 \(\mathcal{C}\) 在常值函数所成之集上诱导的拓扑是离散拓扑，因而 \(\mathcal{C}\) 不与 \(\mathbf{R}\) -向量空间结构（ \(\mathbf{A}\) 的）相容。
\item
  证明 A 关于拓扑 \(\mathcal{C}\) 是一个盖尔范德环（第 III 章，§ 6，习题 11）。
\end{enumerate}

¶ 15) 设 X 是完全正则空间，βX 是它的斯通–切赫紧化。给环 A = C(X; R) 赋以习题 14 中定义的拓扑 C。对每个 f ∈ A，用 V(f) 表示 X 的子集 \(\overline{f}^{-1}(\mathrm{o})\) 在 βX 中的闭包。对 A 的每个子集 M，用 V(M) 表示 ∩f ∈ M V(f)。对 βX 的每个子集 Y，用 ℑ(Y) 表示使 Y ⊂ V(f) 的一切 f ∈ A 所成之集，且用 ℑ(x) 代替 ℑ(\{x\})。我们有 V(f) = ∅，当且仅当 f 是 A 中的单位。

\begin{enumerate}
\def\labelenumi{\alph{enumi})}
\item
  若 \(f, g\) 是 \(\mathbf{A}\) 中使 \(\overline{f}^{-1}(\mathrm{o}) \cap \overline{g}^{-1}(\mathrm{o}) = \emptyset\) 的两个元素，证明 \(\mathrm{V}(f) \cap \mathrm{V}(g) = \emptyset\) 。{[}考察函数 \(|f| / (|f| + |g|)\) 。由此推出：对每个子集 \(\mathbf{M}\) （ \(\mathbf{A}\) 的），我们有 \(\mathrm{V}(\mathbf{M}) = \mathrm{V}(\mathfrak{a})\) ，其中 \(\mathfrak{a}\) 是 \(\mathbf{M}\) 在 \(\mathbf{A}'\) 中生成的理想。
\item
  证明：对 \(\beta X\) 的每个子集 Y，\(\Im(Y)\) 是 A 的理想{[}利用 a){]}，且 \(\mathrm{V}(\Im(\mathrm{Y}))\) 是闭包 \(\overline{Y}\) （Y 在 \(\beta X\) 中的）。
\item
  设 \(u \in \mathbf{A}\) 使 \(u(x) > 0\) 对一切 \(x \in \mathbf{X}\) 成立，并设 \(g \in \mathbf{A}\) 。证明存在 \(f \in \mathbf{A}\) ，使得 \(|f - g| \leqslant u\) 且使得 \(\mathbf{V}(f)\) 是 \(\mathbf{V}(g)\) 的一个邻域。{[}定义 \(f\) 等于 \(g + u\) （当 \(g(x) + u(x) < 0\) 时），等于 \(g - u\) （当 \(g(x) - u(x) > 0\) 时），否则等于 \(o\) ；考察函数 \(f' = \inf ((u + g)^+, (u - g)^-)\) 并利用 \(a)\) 。{]}
\item
  设 \(\mathfrak{a}\) 是 \(\mathbf{A}\) 的理想。证明：若 \(f\in\mathbf{A}\) 使 \(\mathbf{V}(f)\) 是 \(\mathbf{V}(\mathfrak{a})\) 的一个邻域，则 \(f\in\mathfrak{a}\) 。{[}首先注意存在 \(g \in \mathfrak{a}\) 使 \(\mathbf{V}(g)\) 含于 \(\mathbf{V}(f)\) 的内部，然后考察函数 \(h\) （定义在 \(\mathbf{X}\) 上）：它等于 \(f(x) / g(x)\) （当 \(f(x) \neq 0\) 时），而等于 \(0\) （当 \(f(x) = 0\) 时）。{]}
\item
  由 c) 与 d) 推出：\(\Im(\mathrm{V}(\mathfrak{a})) = \overline{\mathfrak{a}}\) 对 A 的每个理想 \(\mathfrak{a}\) 成立。{[}首先证明 \(\mathrm{V}(\mathbf{M}) = \mathrm{V}(\overline{\mathbf{M}})\) 对 A 的所有子集 M 成立，用反证法论证。{]}由此推出：对 \(\beta\mathrm{X}\) 的每个子集 Y，\(\Im(\mathrm{Y})\) 是 A 的闭理想。{[}注意 \(\Im(\mathrm{Y}) = \Im(\mathrm{V}(\Im(\mathrm{Y})))\) 。{]}
\item
  利用 e) 证明：\(x \to \Im(x)\) 是 \(\beta X\) 到 A 的（必然是闭的）极大理想所成之集上的双射。A 的一个理想是闭的，当且仅当它是诸极大理想的一个交。环 A 无根。
\end{enumerate}

¶ 16) 我们保留习题 15 的假设与记号。

\begin{enumerate}
\def\labelenumi{\alph{enumi})}
\item
  设 a 是环 A 中的闭理想。证明：在商环 A/a 中，A 中函数 \(f \geqslant 0\) 的典范像所成之集 P，是关于某个使 A/a 成为格序环的序的、元素 \(\geqslant 0\) 所成之集{[}利用习题 15 e) 注意：若 f 与 g 属于 a，则 \(|f| + |g|\) 以及每个函数 \(h \in A\) （它使 \(|h| \leqslant |f|\) ）也属于 a。常值函数所成之集在 A/a 中的典范像作为一个有序域同构于 R。
\item
  特别地，对每个 \(x \in \beta X\) ，域 \(\mathrm{A}/\Im(x)\) 典范地赋有一个有序域结构。称极大理想 \(\Im(x)\) 为实的，如果 \(\mathrm{A}/\Im(x)\) 同构于 R；否则称为超实的。\(\Im(x)\) 是超实的，其充分必要条件是：\(x \in \beta X - X\) 且存在单位 \(f \in A\) 使 \(\lim_{y \to x, y \in X} f(y) = 0\) 。由此推出：A 的所有极大理想都是实的，当且仅当 X 是伪紧的（第 IX 章，§1，习题 21）。
\item
  对每个 \(x \in \beta\mathbf{X}\) ，设 \(\varphi_x\) 是典范同态 \(\mathbf{A} \to \mathbf{A}/\Im(x)\) 。为使 \(\varphi_x(f)\) 是 \(\geqslant 0\) 的（在 \(\mathbf{A}/\Im(x)\) 中），必须且只需存在 \(g \in \Im(x)\) 使 \(f(y) \geqslant 0\) 在 \(\overline{g}^{-1}(0)\) 中成立{[}注意关系 \(\varphi_x(f) \geqslant 0\) 等价于 \(f \equiv |f|\pmod{\Im(x)}\) {]}。为使 \(\varphi_x(f)\) 是 \(>0\) 的，必须且只需存在 \(g \in \Im(x)\) 使 \(f(y) > 0\) 在 \(\overline{g}^{-1}(0)\) 中成立。{[}注意：若 \(f \notin \Im(x)\) ，则存在 \(g \in \Im(x)\) 使 \(\overline{f}^{-1}(0) \cap \overline{g}^{-1}(0) = \emptyset\) （利用习题 15）。{]}
\item
  证明：若 \(\Im(x)\) 是超实的，则域 \(\mathbf{A}/\Im(x)\) 在 \(\mathbf{R}\) 上的超越次数至少等于 \(\mathfrak{c}=\operatorname{Card}(\mathbf{R})\) 。{[}设 \(f>0\) 是 \(\mathbf{A}\) 的单位，使 \(\varphi_x(f)=u\) 关于 \(\mathbf{R}\) 是无穷大。证明：当 \(r\) 跑遍 \(\mathbf{R}\) 在 \(\mathbf{Q}\) 上的一个基（由数 \(>0\) 组成）中的元素时，元素 \(\varphi_x(f^r)\) （ \(\mathbf{A}/\Im(x)\) 的）是代数无关的。{]}
\end{enumerate}

* e) 对每个点 \(a = (a_{1}, \ldots, a_{n}) \in \mathbb{R}^{n}\) ，用 \(f_{a}(\mathbf{X})\) 表示多项式 \(\mathbf{X}^{n} + a_{1}\mathbf{X}^{n-1} + \cdots + a_{n}\) 。对每个实数 \(\xi\) ，令 \(\nu(\xi)\) 为零点 \(z\) （ \(f_{a}\) 在 \(\mathbf{C}\) 中的）中使 \(\mathrm{R}(z) = \xi\) 者的重数之和；又对每个整数 \(k\) （它使 \(\mathrm{i} \leqslant k \leqslant n\) ），令 \(\rho_{k}(a)\) 为最小的实数 \(\xi\) ，使得 \(\sum_{\eta \leqslant \xi} \nu(\eta) \geqslant k\) 。证明每个函数 \(\rho_{k}\) 都在 \(\mathbb{R}^{n}\) 上连续（利用鲁歇定理）。

\begin{enumerate}
\def\labelenumi{\alph{enumi})}
\setcounter{enumi}{5}
\item
  证明：对每个 \(x \in \beta X\) ，有序域 \(A/\Im(x)\) 是实闭的。{[}设 \(F(X) = X^{n} + f_{1}X^{n-1} + \cdots + f_{n}\) 是系数在 A 中的奇数次多项式；函数 \(g_{k}(y) = \rho_{k}(f_{1}(y), \ldots, f_{n}(y))\) （其中 \(\rho_{k}\) 是 e) 中定义的函数）在 X 上连续，且对每个 \(y \in X\) 存在指标 k 使 \(F(g_{k}(y)) = 0\) ；由此推出乘积 \(F(g_{1}) \cdots F(g_{n})\) 属于 \(\Im(x).]\) *
\item
  设 X 是无限离散空间，\(y \to F_y\) 是 X 到 X 的有限子集所成之集上的一个双射；对每个 \(z \in X\) ，令 \(M_z\) 为一切 \(y \in X\) 中使 \(z \in F_y\) 者所成之集；这些集构成 X 上一个滤子 F 的基。设 x 是 \(\beta X\) 中的点，使 X 上对应的超滤子（第 I 章，§9，习题 27）比 F 细。设 B 是 A 的子集，使 Card (B) ≤ Card (X)，并设 \(y \to g_y\) 是满射 \(X \to B\) 。对每个 \(z \in X\) ，令 \(f(z) = \mathrm{i} + \sup_{y \in \mathbb{F}_z} g_y(z)\) 。证明 \(f(z) > g_y(z)\) 对一切 \(z \in M_y\) 成立，并由此推出 \(\varphi_x(f) > \varphi_x(g_y)\) 对一切 \(y \in Y\) 成立{[}利用 c){]}。由此得出结论：Card \(A/\Im(x)\) \textgreater{} Card X。
\end{enumerate}

¶ 17) 我们保留习题 15 的假设与记号。

\begin{enumerate}
\def\labelenumi{\alph{enumi})}
\item
  极大理想 \(\Im(x)\) 是实的，当且仅当：对每个无限序列 \((g_n)\) （由 \(\Im(x)\) 的元素组成），诸集 \(\overline{g_n}^{-1}(0)\) 的交非空。{[}利用习题 16 b)；为证条件是必要的，证明：若 \(\bigcap_{n}\overline{g_{n}}^{-1}(0) = \emptyset\) ，则函数 \(f(y) = \sum_{n = 0}^{\infty}\inf (g_n(y),2^{-n})\) 连续，是 A 中的单位，且当 \(y\) 保持于 X 中而趋于 \(x\) 时趋于 o；注意 \(f(y)\leqslant 2^{-m}\) （对 \(y\in \bigcap_{k\leqslant m}\overline{g_k}^{-1}(0)\) ），并利用习题 15 a) 以及下述事实：诸 V(g)（其中 \(g\in \Im(x)\) ）构成 \(x\) 在 \(\beta \mathbf{X}.\) 中的邻域基本系。{]}
\item
  称 X 是实紧的，如果仅有的实极大理想 \(\Im(x)\) 是那些使 \(x \in \mathbf{X}\) 的理想。证明每个完全正则的林德勒夫空间（第 I 章，§ 9，习题 15）都是实紧的{[}利用 \(a\) {]}。一个完全正则的伪紧空间，除非它是紧的，否则不是实紧的。
\item
  设 \(\nu\mathbf{X}\) 是一切 \(x\in \beta \mathbf{X}\) 中使 \(\Im (x)\) 为实者所成之集。证明每个连续实值函数 \(f\in \mathbf{A}\) 都可连续扩张到 \(\nu\mathbf{X}\) ，从而 \(\mathcal{C}(\mathbf{X};\mathbf{R})\) 与 \(\mathcal{C}(\upsilon\mathbf{X};\mathbf{R})\) 可典范地等同。证明子空间 \(\nu\mathbf{X}\) （ \(\beta\mathbf{X}\) 的）是实紧的{[}注意：若 \(f\) 是 \(\mathcal{C}(\mathbf{X};\mathbf{R})\) 中的单位，则它到 \(\nu\mathbf{X}\) 上的连续扩张是 \(\mathcal{C}(\upsilon\mathbf{X};\mathbf{R})\) 中的单位，且 \(\beta(\upsilon\mathbf{X}) = \beta\mathbf{X}]\) 。\(\nu\mathbf{X}\) 称为 \(\mathbf{X}\) 的实紧化。
\item
  证明 \(\nu\mathbf{X}\) 是 \(\mathbf{X}\) 关于 \(\mathbf{X}\) 上使一切函数 \(f\in \mathbf{A}\) 都一致连续的最粗一致结构的完备化。（注意柯西滤子 \(\mathfrak{F}\) （ \(\mathbf{X}\) 上的，关于这个一致结构）收敛于一点 \(x\in \beta \mathbf{X}\) ；若 \(x\notin \mathfrak{vX}\) ，则将有函数 \(f\in \mathbf{A}\) ，它在 F 的某个集合上无界。{]}
\item
  由 d) 推出：若 \(u: X \to Y\) 是完全正则空间 X 到实紧空间 Y 中的任一连续映射，则 u 可连续扩张为映射 \(\upsilon X \to Y\) 。于是空间 \(\nu X\) 是如下泛映射问题的解：其中 \(\Sigma\) 是实紧空间结构，而 \(\alpha\) -映射与态射是连续映射（《集合论》，第 IV 章，§ 3，第 1 号）。
\end{enumerate}

\(f\) ) 证明实紧空间的每个闭子空间是实紧的，实紧空间的每个乘积是实紧的，且在豪斯多夫空间中，实紧子空间的任意交是实紧的。{[}例如，为证闭子空间 \(\mathbf{X}\) （实紧空间 \(\mathbf{Y}\) 的）是实紧的，考察扩张 \(\vec{j}\) （到 \(\nu\mathbf{X}\) 上的，它来自典范单射 \(j:\mathbf{X}\to\mathbf{Y}\) ），并注意 \(j^{-1}(\mathbf{X})=\mathbf{X}\) ；由此推出 \(\mathbf{X}\) 在 \(\nu\mathbf{X}\) 中是闭的。另两种情形的证明类似。{]}

\begin{enumerate}
\def\labelenumi{\alph{enumi})}
\setcounter{enumi}{6}
\tightlist
\item
  由 b)、d) 与 f) 推出：一个完全正则空间是实紧的，当且仅当它同胚于乘积 \(R^{I}\) 的一个闭子空间。
\end{enumerate}

\begin{enumerate}
\def\labelenumi{\arabic{enumi})}
\setcounter{enumi}{17}
\tightlist
\item
  设 X 是实紧空间。
\end{enumerate}

\begin{enumerate}
\def\labelenumi{\alph{enumi})}
\item
  证明：对每个非零同态 \(u: \mathcal{C}(\mathbf{X}; \mathbf{R}) \to \mathbf{R}\) ，存在唯一的点 \(x \in \mathbf{X}\) ，使得 \(u(f) = f(x)\) 对一切 \(f\) （ \(\mathcal{C}(\mathbf{X}; \mathbf{R})\) 中的）成立。
\item
  设 Y 是另一个实紧空间。证明：对每个把单位元映为单位元的 R-代数同态 \(v: \mathcal{C}(Y; \mathbf{R}) \to \mathcal{C}(X; \mathbf{R})\) ，存在唯一的连续映射 \(w: X \to Y\) ，使得 \(v(g) = g \circ w\) 对一切 \(g \in \mathcal{C}(Y; \mathbf{R})\) 成立。特别地，\(\mathcal{C}(X; \mathbf{R})\) 与 \(\mathcal{C}(Y; \mathbf{R})\) 同构，当且仅当 X 与 Y 同胚。
\end{enumerate}

¶ 19) 设 X 是紧空间，A 是赋范 C-代数 C(X; C)。对 A 的每个闭理想 a，证明关系 \(f \in a\) 蕴含 \(\overline{f} \in a\) （注意 \(\overline{f}\) 是形如 \(g f\) 的函数的一致极限）。把习题 13 的结果推广到代数 A。

¶ 20) a) 设 X 是紧空间，A 是有单位元的有限秩 R-代数，并给 A 赋以第 VI 章，§ 1，第 5 号中定义的拓扑。设 B 是 R-代数 C (X; A) 的子代数，且假设：对每个 ε \textgreater{} o、X 的每对不同点 x, y 以及 A 的每对元素 u, v，存在 f ∈ B 使 \textbar\textbar f(x) - u\textbar\textbar{} ≤ ε 且 \textbar\textbar f(y) - v\textbar\textbar{} ≤ ε （ \textbar\textbar{} \textbar\textbar{} 为定义 A 的拓扑的任一范数）；进而假设：当 u 与 v 属于 R 时，存在满足这些条件且取值于 R 中的函数 f ∈ B。证明在这些假设下，B 关于一致收敛拓扑在 C (X; A) 中稠密。{[}首先证明 C (X; R) 的每个函数都可用 B 中函数一致逼近，然后利用假设和一个单位分解证明 A 中每个常元 a ∈ A 也是如此。{]}

\begin{enumerate}
\def\labelenumi{\alph{enumi})}
\setcounter{enumi}{1}
\tightlist
\item
  取 A = H，即 R 上的四元数体。设 B 是 R-代数 \(\mathcal{C}(\mathbf{X}; \mathbf{H})\) 的子代数，满足：(i) 对每个 \(f \in B\) ，f 的共轭 \(\overline{f}\) {[}由 \(\overline{f}(x) = \overline{f(x)}\) （对所有 \(x \in X\) ）定义{]}属于 B；(ii) 对每个 \(\varepsilon > 0\) 、X 的每对不同点 x, y 以及 H 的每对元素 u, v，存在 \(f \in B\) 使 \(||f(x) - u|| \leqslant \varepsilon\) 且 \(||f(y) - v|| \leqslant \varepsilon\) 。证明 B 在 \(\mathcal{C}(\mathbf{X}; \mathbf{H})\) 中稠密{[}利用 a){]}。c) 把习题 13 的结果推广到代数 \(\mathcal{C}(\mathbf{X}; \mathbf{H})\) ，并特别证明这个代数的每个闭理想都是双侧的{[}利用 b)，按习题 19 的方式论证{]}。
\end{enumerate}

¶ 21) a) 设 X 是全不连通紧空间，A 是拓扑环。给环 C (X; A) 赋以一致收敛拓扑，它与环结构相容。设 a 是 C (X; A) 的左理想，且对每个 \(x \in X\) ，设 \(\Im(x)\) 是诸 \(f(x)\) （ \(f \in a\) ）所成之集在 A 中的闭包；\(\Im(x)\) 是 A 的闭左理想。证明 \(\overline{a}\) 恒同于一切 \(f \in C(X; A)\) 中使 \(f(x) \in \Im(x)\) （对所有 \(x \in X\) ）成立者所成之集（注意对每个 \(x \in X\) ，x 的既开又闭的邻域构成 x 的邻域基本系）。

\begin{enumerate}
\def\labelenumi{\alph{enumi})}
\setcounter{enumi}{1}
\tightlist
\item
  再假设 A 有一个由双侧理想组成的 o 的邻域基本系。设 B 是 \(\mathcal{C}(\mathbf{X};\mathbf{A})\) 的子环，它包含常值函数且分离 X 的点。证明 B 在 \(\mathcal{C}(\mathbf{X};\mathbf{A})\) 中稠密。{[}对 X 的每个闭子集 F、每个 \(x\notin F\) 以及 A 中 o 的每个邻域 U，证明存在 \(f\in B\) 使 \(f(x)=1\) 且 \(f(y)\in U\) 对一切 \(y\in F\) 成立。{]}
\end{enumerate}

\section*{历史注记}
\phantomsection
\addcontentsline{toc}{section}{历史注记}

（方括号中的数字指本注记末尾的参考文献。）

在十九世纪初，任意函数的概念几乎无人知晓。更不用说，一般地研究函数的集合并给它们赋予拓扑结构的思想，在黎曼的时代之前并未出现（见第 I 章的历史注记），而且直到十九世纪末，这一思想才得到系统的运用。

然而，实值函数序列收敛的思想，自从无穷小分析诞生以来就或多或少地被自觉不自觉地使用着。当然，这里所谓收敛指的是逐点收敛；其他的收敛类型，要等到收敛级数与连续函数的概念由波尔查诺和柯西精确加以定义之后，才可能被描述。后者起初并未认识到逐点收敛与一致收敛之间的区别，并且相信自己已经证明，任何连续函数的收敛级数之和都是连续的（{[}1{]}, (2), vol.~3, p.~120）。这个错误几乎立即就被阿贝尔指出，他同时用一个经典的论证表明，每个幂级数在其收敛区间内部都是连续的；这个论证在上述特殊情形中本质上使用了一致收敛的概念（{[}2{]}, pp.~223-224）。剩下的事只是把这一概念一般地表述出来，这项工作由斯托克斯和赛德尔于 1847-1848 年独立完成，柯西本人也于 1853 年完成（{[}1{]}, (1), vol.~12, p.~30）(*)。

在魏尔斯特拉斯与黎曼的影响下，一致收敛概念及有关问题的系统研究，在十九世纪最后三分之一的时间里由德国学派（汉克尔、杜·布瓦–雷蒙）、而尤其是由意大利学派发展起来；狄尼与阿尔泽拉明确了连续函数序列的极限为连续的必要条件，而阿斯科利则引入了等度连续这一基本概念，并证明了刻画连续函数紧集的定理{[}3{]}（这个定理后来由蒙泰尔在其“正规族”理论中加以普及，正规族即解析函数的相对紧集）。

另一方面，魏尔斯特拉斯本人发现了 \([4b]\) 在有界集上用多项式一致逼近一个或多个实变量的连续实值函数的可能性。这一结果立即引起了浓厚的兴趣，并引出了许多“定量的”研究(*)。

对这些问题，现代的贡献主要在于赋予它们以所能达到的完全一般性：考虑定义域与值域不再限于 R 或有限维空间的函数，从而借助一般拓扑概念把它们安置在其固有的环境中。特别地，魏尔斯特拉斯定理在经典分析中早已显示出是一件强有力的武器，近年来 M. H. 斯通把它推广到了广泛得多的场合；他发展了 H. 勒贝格引进的一个思想（在魏尔斯特拉斯定理的一个证明中），清楚地指明了函数格在连续实值函数逼近理论中所起的重要作用（用“格多项式”逼近，参看 § 4 的命题 2 与定理 2），并且表明了推广的魏尔斯特拉斯定理如何以一整系列类似的逼近定理作为其直接推论，这些定理因而能够以连贯得多的方式归拢在一起。我们或多或少遵循了他的阐述{[}5{]}。

\subsection*{参考文献}

{[}1{]} A.-L. CAUCHY, Œuvres, Paris (Gauthier-Villars), 1882-1932.

{[}2{]} N. H. ABEL, Œuvres, vol.~1, p.~223-224, edited by Sylow and Lie (Christiania) 1881.

{[}3{]} G. Ascoli, Sulle curve limiti di una varietà data di curve, Mem. Accad. Lincei (3), vol.~18 (1883), pp.~521-586.

{[}4{]} K. WEIERSTRASS, Mathematische Werke, Berlin (Mayer und Müller), 1894-1903:

\begin{enumerate}
\def\labelenumi{\alph{enumi})}
\item
  Zur Theorie der Potenzreihen, Bd. 1, pp.~67-74;
\item
  Über die analytische Darstellbarkeit sogenannter willkürlicher Functionen reeller Argumente, Bd. 3, pp.~1-37.
\end{enumerate}

{[}5{]} M. H. STONE, The generalized Weierstrass approximation theorem, Mathematics Magazine, 21 (1948), pp.~167-183 and 237-254.

{[}6{]} C. DE LA VALLÉE-POUSSIN, Leçons sur l'approximation des fonctions d'une variable réelle, Paris, (Gauthier-Villars), 1919.

引用数字依次表示章、节和小节或习题。

T : V, I, 2.

\(a^{x}\) （a 为任意实数 \textgreater0，x 为任意实数） : V, 4, I.

\(\log_{a}x\) （a、x 为实数 \textgreater0，\(a \neq i\) ） : V, 4, I.

\(R^{n}\) : VI, I, I.

\(d(x, y)\) （欧氏距离） : VI, 2, I.

\(||x||\) （欧氏范数） : VI, 2, I.

\((x|y)\) （内积） : VI, 2, 2.

\(S_{n}, B_{n}, R_{n}^{*}\) : VI, 2, 3.

\(P_{n}(R), P_{n}\) : VI, 3, I.

\(\tilde{R}\) : VI, 3, 3.

\(\infty\) （ \(\tilde{R}\) 的点） : VI, 3, 3.

\(P_{n,p}(R), P_{n,p}\) : VI, 3, 5.

\(r(G), d(G)\) （G 为 \(R^{n}\) 的任一闭子群） : VII, I, 2.

\(G^{*}\) （G 为 \(R^{n}\) 的任一子群） : VII, I, 3.

\(T^{n}\) : VII, I, 4.

C, i: VIII, I, I.

\(\mathcal{R}(z), \Im(z), \bar{z}, |z|\) （z 为复数）: VIII, I, I.

\(C^{*}, U\) : VIII, I, 3.

H : VIII, I, 4.

\(e(x) (= e^{2\pi ix})\) : VIII, 2, I.

\((\widehat{\Delta_{1}, \Delta_{2}})(\Delta_{1}, \Delta_{2}\text{ 射线})\) : VIII, 2, 2.

\(\delta, \varpi\) : VIII, 2, 2.

Am (z) （z 为复数）: VIII, 2, 2.

cos θ, sin θ, tan θ （θ 为角）: VIII, 2, 2.

\(\cos_{a}x\) , \(\sin_{a}x\) , \(\tan_{a}x\) , \(\cot_{a}x\) （x 为实数）: VIII, 2, 4.

cos x, sin x, tan x, cot x （x 为实数）: VIII, 2, 4.

\((\widehat{D_{1}, D_{2}})(D_{1}, D_{2}\text{ 直线})\) : VIII, 2, 6.

\(\delta_{0}\) : VIII, 2, 6.

\(C^{n}\) : VIII, 4, I.

\(P_{n}(C), \tilde{C}\) : VIII, 4, 3.

∞ （ \(\tilde{C}\) 的点） : VIII, 4, 3.

\[
\mathbf {P} _ {n, p} (\mathbf {C}): \text { VIII }, 4, 4.
\]

\(\beta \mathbf{X}\) （X 为完全正则空间）: IX, 1, 习题 7.

\(d(\mathbf{A}, \mathbf{B}), d(x, \mathbf{A})\) {[} \(x\) 为一点，A、B 为一个度量空间的子集，其中两点间的距离记作 \(d(x, y)] : \mathrm{IX}, 2, 2\)

\(|x|\) （赋值除环中的绝对值）: IX, 3, 2.

\textbar\textbar x\textbar\textbar{} （赋范空间中的范数） : IX, 3, 3.

D(A) : IX, 5, 习题 3.

L(C) （C 为筛）: IX, 6, 5.

\(\prod_{n \in \mathbb{N}} x_n : \text{IX, Appendix, I.}\)

\(P_{n=h} x_n : \text{IX, Appendix, 3.}\)

\(\prod_{n=0}^{n} x_n: \text{IX, Appendix, 3.}\)

\(\mathcal{F}(\mathrm{X};\mathrm{Y}),\mathrm{H}(x)\) {[}H 为 \(\mathcal{F}(\mathrm{X};\mathrm{Y})]\) 的子集，\(\Phi (x)\) {[}\(\Phi\) 为 \(\mathcal{F}(\mathrm{X};\mathrm{Y})]\) 上的一个滤子，\(u|\mathrm{A},\mathrm{H}|\mathrm{A}\) {[}A 为 \(\mathrm{X},u\in \mathcal{F}(\mathrm{X};\mathrm{Y}),\mathrm{H}\subset \mathcal{F}(\mathrm{X};\mathrm{Y})]:\mathrm{X},\) 1.

\(W(\mathbf{V})\) （V 为 \(\mathbf{Y}\) 的一个围域） : \(\mathbf{X}, \mathrm{I}, \mathrm{I}\) .

\[
\mathcal {F} _ {u} (\mathrm{X}; \mathrm{Y}): \mathrm{X}, \mathrm{I}, \mathrm{I}.
\]

\(\mathcal{F}_{\mathfrak{S}}(\mathbf{X};\mathbf{Y}):\mathbf{X},\mathbf{1},\mathbf{2}.\)

\(W(A, V) (A \in \mathfrak{S}, V \text{ 是一个围域，相对于 } Y): X, I, 2.\)

\[
\mathcal {F} _ {s} (\mathrm{X}; \mathrm{Y}), \mathcal {F} _ {c} (\mathrm{X}; \mathrm{Y}): \mathrm{X}, \mathrm{I}, 3.
\]

\[
\mathcal {C} (\mathbf {X}; \mathbf {Y}), \mathcal {C} _ {\mathfrak {S}} (\mathbf {X}; \mathbf {Y}), \mathcal {C} _ {s} (\mathbf {X}; \mathbf {Y}), \mathcal {C} _ {c} (\mathbf {X}; \mathbf {Y}), \mathcal {C} _ {u} (\mathbf {X}; \mathbf {Y}): \mathbf {X}, \mathrm{I}, 6.
\]

\[
\tilde {\mathcal {C}} _ {\mathfrak {S}} (\mathrm{X}; \mathrm{Y}): \mathrm{X}, \mathrm{I}, 6.
\]

\(\tilde{x}\) （x 为 X 的点）: X, 2, 1.

\(\mathcal{B}_{\mathfrak{S}}(\mathbf{X};\mathbf{Y}),\mathcal{B}(\mathbf{X};\mathbf{Y})\) （Y 为度量空间）: \(\mathbf{X},3,\mathrm{i}\)

\(||u||\) （u 为到赋范空间中的有界映射）: X, 3, 2.

\(\mathscr{L}\left(\mathbf{X}_{1},\ldots,\mathbf{X}_{n};\mathbf{Y}\right)\left(\mathbf{X}_{1},\ldots,\mathbf{X}_{n}\text{ 和 Y 均为赋范空间}\right):\mathbf{X},3,2.\)

\(||u||\) （u 为赋范空间的乘积到赋范空间中的多重线性映射）: X, 3, 2.

\(T(\mathbf{K},\mathbf{U})\) （K 为 \(\mathbf{X}\) 的紧子集，U 为 \(\mathbf{Y}\) 的开子集） : \(\mathbf{X}\) , 3, 4. \(\mathcal{C}_{\beta}:\mathbf{X}\) , 3, 5.

\(\mathcal{C}^{\infty}\) (X; R): X, 4, 习题 6.

\(\nu X:X,4,\) 习题 17.

引用数字依次表示章、节和小节或习题。

横坐标 : VI, 1, 4. 横轴 : VI, 1, 4. 绝对值 : IX, 3, 2. 平凡绝对值 : IX, 3, 2. p 进绝对值 : IX, 3, 2. 等价绝对值 : IX, 3, 2. 复数的绝对值 : VIII, 1, 1. 复数的绝对收敛无穷乘积 : VIII, 3, 3. 绝对收敛级数（在赋范空间中）: IX, 3, 6. R\^{}n 中点的绝对收敛级数 : VII, 3, 2. 绝对可和族（在赋范空间中）: IX, 3, 6. 锐角域 : VIII, 2, 5. R\^{}n 的加法群 : VI, 1, 2. 模 a 的实数加法群 : V, 1, 2. R\^{}n 的加法一致结构 : VII, 1, 2. 容许子集 : IX, 6, 习题 12. 相互垂直的仿射线性簇 : VI, 2, 2. 赋范代数 : IX, 3, 7. 复数的代数范数 : VIII, 1, 1. 代数数 : VIII, 1, 习题 2. 几乎开映射 : IX, 5, 习题 24. 几乎开集 : IX, 5, 习题 6. 概周期函数 : X, 3, 习题 28. 概周期点 : X, 3, 习题 27. 复数的辐角 : VIII, 2, 2 与 VIII, 2, 3. 平角 : VIII, 2, 2. 一对射线的角 : VIII, 2, 2 与 VIII, 2, 3. 角域的角 : VIII, 2, 5. 正直角 : VIII, 2, 2.

角域（锐的、直的、钝的、平的、凸的、凹的）: VIII, 2, 5. 一致逼近: X, 4, 1. 阿基米德的（线性有序群）: V, 3, 习题 1. 阿基米德公理: V, 2. 复数的幅角: VIII, 2, 2. 阿斯科利定理: X, 2, 5. \(\mathbf{R}^n\) 的子群的相伴子群 : VII, 1, 3. 阿基米德公理: V, 2. 坐标轴: VI, 1, 4. 虚轴: VIII, 1, 2. 横轴: VI, 1, 4. 纵轴: VI, 1, 4. 实轴: VIII, 1, 2.

紧收敛的一致结构 : X, 1, 3.

紧开拓扑 : X, 3, 4.

斯通–切赫紧化 : IX, 1, 习题 7.

相容的（度量与拓扑） : IX, 2, 5.

相容的（度量与一致结构） : IX, 2, 4.

相容的（范数与代数结构） : IX, 3, 7.

完全正规空间 : IX, 4, 习题 3.

完全正则滤子 : IX, 1, 习题 8.

完全正则空间 : IX, 1, 5.

与可一致化空间相伴的完全正则空间 : IX, 1, 习题 4.

完全分离的（完全正则空间中的闭集）: IX, 1, 习题 11.

\(\mathbf{C}^n\) 中的复超平面 : VIII, 4, 1.

\(\mathbf{C}^n\) 中的复线性簇 : VIII, 4, 1.

\(\mathbf{C}^n\) 中的复直线 : VIII, 4, 1.

复数 : VIII, 1, 1.

n 维复数空间 : VIII, 4, 1.

\(\mathbf{C}^n\) 中的复平面 : VIII, 4, 1.

复射影直线 : VIII, 4, 3.

复射影平面 : VIII, 4, 3.

n 维复射影空间 : VIII, 4, 3.

复数的共轭 : VIII, 1, 1.

连续单位分解 : IX, 4, 3.

收敛的无穷乘积（在赋范代数中） : IX, 附录, 3.

正规收敛的 : X, 3, 2.

逐点收敛的 : X, 1, 3.

一致收敛的 : X, 1, 1.

坐标轴 : VI, 1, 4.

坐标簇 : VI, 1, 4.

角的余弦 : VIII, 2, 2.

数的余弦 : VIII, 2, 4.

角的余切 : VIII, 2, 2.

数的余切 : VIII, 2, 4.

可数型（可度量化空间） : IX, 2, 8.

可数紧空间 : IX, 2, 习题 14.

可除覆盖 : IX, 4, 习题 16.

均匀覆盖 : IX, 4, 习题 16.

逐点有限覆盖 : IX, 4, 3.

一对直线的叉角 : VIII, 2, 6.

直叉角 : VIII, 2, 6.

方体 : IX, 1, 5.

立方体（开的、闭的） : VI, I, I.

佩亚诺曲线 : VI, 1, 习题 2.

度（角度量的单位） : VIII, 2, 3.

直径 : IX, 2, 2.

直径超平面 : VI, 2, 4.

\(\mathbf{S}_n:\mathbf{VI},3,\mathbf{i}\) 的对径点

\(\mathbf{R}^n\) 的闭子群的维数 : VII, 1, 2.

迪尼定理 : X, 4, 1.

直线或射线的方向比 : VI, 1, 4.

直线或射线的方向向量 : VI, 1, 4.

圆盘（开的、闭的） : VI, 2, 3.

子集的离散族 : IX, 4, 习题 18.

欧氏位移 : VI, 2, 2.

两集合间的距离 : IX, 2, 2.

欧氏距离 : VI, 2, 1.

点到集合的距离 : IX, 2, 2.

可除覆盖 : IX, 4, 习题 16.

四元数体 : VIII, 1, 4.

赋值除环 : IX, 3, 2.

\(\varepsilon\) -周期（函数的）: X, 3, 习题 28.

等度连续的映射集 : X, 2, 1.

在一点等度连续 : X, 2, 1.

一致等度连续 : X, 2, 1.

模 i 均分 : VII, i, 习题 14.

等价绝对值 : IX, 3, 2.

等价的伪度量族 : IX, 1, 2.

等价范数 : IX, 3, 3.

等价伪度量 : IX, 1, 2.

等价半绝对值 : IX, 3, 习题 10.

映入 \(\mathbf{P}_n\) 的本质映射 : VI, 3, 习题 2.

映入 \(\mathbf{S}_1\) 的本质映射 : VI, 2, 习题 6.

球（闭的、开的） : VI, 2, 3.

欧氏位移 : VI, 2, 2.

欧氏距离 : VI, 2, 1.

\(\mathbf{R}^n\) 中的欧氏范数 : VI, 2, 1.

球面 : VI, 2, 3.

均匀覆盖 : IX, 4, 习题 16.

指数函数 : V, 4, 1.

乘积的因子（在赋范代数中） : IX, 附录, 1.

绝对可和族 : IX, 3, 6.

从属于子集族的函数族 : IX, 4, 3.

法里数列 : VII, 1, 习题 13.

复数域 : VIII, 1, 1.

完全正则滤子 : IX, 1, 习题 8.

极大完全正则滤子 : IX, 1, 习题 8.

有限开覆盖的一致结构 : IX, 4, 习题 17.

平角 : VIII, 2, 2.

平角域 : VIII, 2, 5.

概周期函数 : IX, 3, 习题 28.

指数函数 : V, 4, 1.

周期函数（定义在 \(\mathbf{R}^n\) 上）: VII, 1, 6.

q 重周期函数 : VII, 1, 6.

由一组子集生成的（σ 代数） : IX, 6, 3.

百分度（角度量的单位） : VIII, 2, 3.

加法群（ \(\mathbf{R}^n\) 的）: VI, 1, 2.

加法群（模 a 的实数的）: V, 1, 2.

可度量化群 : IX, 3, 1.

欧氏位移群 : VI, 2, 2.

单参数群 : V, 3.

正交群 : VI, 2, 2.

半直线 : VI, 1, 4.

由超平面确定的半空间（开的、闭的） : VI, 1, 4.

豪斯多夫半绝对值 : IX, 3, 习题 9.

与半绝对值相伴的豪斯多夫半绝对值 :

半球面（开的、闭的） : VI, 2, 4.

分式线性函数 : VI, 3, 5.

超实极大理想 : X, 4, 习题 16.

无穷远超平面 : VI, 3, 3.

复超平面（在 \(\mathbf{C}^n\) 中）: VIII, 4, 1.

直径超平面 : VI, 2, 4.

投影超平面（在球极投影中） : VI, 2, 4.

虚轴 : VIII, 1, 2.

\(\mathbf{R}^n\) 中的格 : VII, I, I.

左不变度量 : IX, 3, 1.

左不变伪度量 : IX, 3, 习题 1.

\(\mathbf{R}^n:\mathrm{VI},2,\mathrm{i}\) 中线段的长度

折线 : VI, 1, 习题 6.

简单折线 : VI, 1, 习题 9.

复射影直线 : VIII, 4, 3.

实射影直线 : VI, 3, 1.

函数的线性组合（系数在赋范空间中）: X, 4, 4.

复线性簇（在 \(\mathbf{C}^n\) 中）: VIII, 4, 1.

线性可达的 : IX, 6, 习题 11.

复直线（在 \(\mathbf{C}^n\) 中）: VIII, 4, 1.

局部仿紧空间 : IX, 4, 习题 27.

以 \(a:\mathbf{V},4,\mathrm{i}\) 为底的对数

对数螺线 : VIII, 3, 习题 5.

卢津空间 : IX, 6, 4.

几乎开映射 : IX, 5, 习题 24.

波莱尔映射 : IX, 6, 习题 16.

由筛分诱导的映射 : IX, 6, 5.

映入 \(\mathbf{P}_n\) 的映射，本质的或非本质的 : VI, 3, 习题 2.

映入 \(S_{1}\) 的映射，本质的或非本质的 : VI, 2, 习题 6.

等距映射 : IX, 2, 2.

分段线性映射 : VI, 1, 习题 6.

极大完全正则滤子 : X, 1, 习题 8.

贫集 : IX, 5, 2.

角的测度 : VIII, 2, 3.

叉角的测度 : VIII, 2, 6.

主测度（角的）: VIII, 2, 3.

主测度（叉角的）: VIII, 2, 6.

亚紧空间 : IX, 4, 习题 25.

度量 : IX, 2, 1.

与伪度量相伴的度量 : IX, 2, 1.

与拓扑相容的度量 : IX, 2, 5.

与一致结构相容的度量 : IX, 2, 4.

左（右）不变度量 : IX, 3, 1.

度量空间 : IX, 2, 1.

可度量化拓扑群 : IX, 3, 1.

可度量化拓扑空间 : IX, 2, 5.

可数型的可度量化拓扑空间 : IX, 2, 8.

可度量化一致空间 : IX, 2, 4.

可度量化一致结构: IX, 2, 4.

可乘序列（在赋范代数中） : IX, 附录, 1.

长田–斯米尔诺夫定理 : IX, 4, 习题 22.

负实半轴 : VIII, 1, 2.

非贫空间 : IX, 5, 习题 7.

范数（在向量空间上）: IX, 3, 3.

代数范数（复数的）: VIII, 1, 1.

与代数结构相容的范数 : IX, 3, 7.

欧氏范数（在 \(\mathbf{R}^n\) 中）: VI, 2, 1.

正规空间 : IX, 4, 1.

正规收敛级数 : X, 3, 2.

赋范代数 : IX, 3, 7.

赋范空间 : IX, 3, 3.

等价范数 : IX, 3, 3.

无处稠密集 : IX, 5, 1.

代数数 : VIII, 1. 习题 2.

复数 : VIII, 1, 1.

n 维复数空间 : VIII, 4, 1.

n 维实数空间 : VI, I, I.

钝角域 : VI, 2, 5.

单参数群 : V, 3.

开球 : IX, 2, 2.

反向射线 : VI, 1, 4.

有序分划 : IX, 附录, 习题 2.

纵坐标 : VI, I, 4.

纵轴 : VI, 1, 4.

射线的端点 : VI, 1, 4.

相互垂直的仿射线性簇 : VI, 2, 2.

正交群 : VI, 2, 2.

正交变换 : VI, 2, 2.

正交向量、正交向量子空间 : VI, 2, 2.

函数在一点的振幅 : IX, 2, 3.

函数在集合中的振幅 : IX, 2, 3.

\(p\) -进绝对值: IX, 3, 7.

超平行体（开的、闭的） : VI, 1, 3.

连续单位分解 : IX, 4, 3.

有序分划 : IX, 附录, 习题 2.

佩亚诺曲线 : VI, 1, 习题 2.

完备正规空间 : IX, 4, 习题 7.

周期函数（定义在 \(\mathbf{R}^n\) 上）: VII, 1, 6.

周期函数的周期 : VII, 1, 6.

分段线性映射 : VI, I, 习题 6.

抽屉原理 : VII, 1, 习题 11.

复射影平面 : VIII, 4, 3.

实射影平面 : VI, 3, 1. 复平面（在 C\(^{n}\) 中） : VIII, 4, 1. 无穷远点 : VI, 3, 3. 逐点有限覆盖 : IX, 4, 3. 一致收敛点 : X, 1, 习题 9. 在 X 的子集上的逐点收敛 : X, 1, 3. 逐点收敛的拓扑 : X, 1, 3 与 3, 4. 逐点收敛的一致结构 : X, 1, 3. 逐点收敛的（滤子） : X, 1, 3. 波兰空间 : IX, 6, 1. 属于给定集合的函数的多项式 : X, 4, 2 与 4, 4. 三角多项式 : X, 4, 4. 正实半轴 : VIII, 1, 2. 预紧收敛的一致结构 : X, 1, 3. 角的主测度 : VIII, 2, 3. 叉角的主测度 : VIII, 2, 6. q 重周期函数的基本周期系 : VII, 1, 6. 绝对收敛的乘积（复数的） : VIII, 3, 3. 无穷乘积 : IX, 附录, 3. 赋范代数中可乘序列的乘积 : IX, 附录, 1. 内积 : VI, 2, 2. 中心投影 : VI, 2, 3. 投影超平面 : VI, 2, 4. 球极投影 : VI, 2, 4. 投影顶点 : VI, 2, 4. n 维复射影空间 : VIII, 4, 3. n 维实射影空间 : VI, 3, 1. 在一点真作用（算子群） : X, 2, 习题 18. 伪紧空间 : IX, 1, 习题 22. 伪度量 : IX, 1, 1. 不变伪度量 : IX, 3, 习题 1. 等价伪度量 : IX, 1, 2. 纯虚数（复数） : VIII, 1, 1.

R\(^{n}\) 上的 q 重周期函数 : VII, 1, 6. 二次超曲面（在 P\(_{n}\) 中）: VI, 3, 习题 10. 二次曲面（在 R\(^{n}\) 中）: VI, 2, 习题 10. 二次锥面（在 P\(_{n}\) 中）: VI, 3, 习题 11. 二次锥面（在 R\(^{n}\) 中）: VI, 2, 习题 11. 四元数 : VIII, 1, 4.

弧度测度 : VIII, 2, 3. 球的半径 : IX, 2, 2. 球或球面的半径 : VI, 2, 3 与 IX, 2, 2.

\(\mathbf{R}^n\) 的子群的有理秩 : VII, 1.

方向比 : VI, 1, 4.

射线（闭的、开的） : VI, 1, 4.

反向射线 : VI, 1, 4.

实轴 : VIII, 1, 2.

实紧空间 : X, 4, 习题 17.

空间的实紧化 : X, 4, 习题 17.

实线性簇（在 \(\mathbf{C}^n\) 中）: VIII, 4, i.

实极大理想: X, 4, 习题 16.

\(n\) 维实数空间 : VI, I, I.

复数的实部 : VIII, 1, 1.

实射影直线 : VI, 3, 1.

实射影平面 : VI, 3, 1.

\(n\) 维实射影空间 : VI, 3, 1.

除环的实赋值 : IX, 3, 2.

矩形（开的、闭的） : VI, I, I.

凹角域 : VIII, 2, 5.

正直角 : VIII, 2, 2.

直角域 : VIII, 2, 5.

直叉角 : VIII, 2, 6.

右不变度量 : IX, 3, 1.

右不变伪度量 : IX, 3, 习题 1.

\(\sigma\) -代数: IX, 6, 3.

凸角域 : VIII, 2, 5.

伪度量的饱和族 : IX, 1, 2.

子集的饱和集 : X, 1, 习题 5.

内积 : VI, 2, 2.

角域 : VIII, 2, 5.

线段（闭的、开的、在 x 处开而在 y 处闭的） : VI, 1, 4.

线段的长度 : VI, 2, 1.

半绝对值 : IX, 3, 习题 9.

豪斯多夫半绝对值 : IX, 3, 习题 9.

等价半绝对值 : IX, 3, 习题 10.

负实半轴 : VIII, 1, 2.

正实半轴 : VIII, 1, 2.

半圆（开的、闭的） : VI, 2, 4.

可分可度量化空间 : IX, 2, 8.

分离的函数集 : X, 4, 1.

集合的点被函数集分离 : X, 4, 1.

可乘序列 : IX, 附录, 1.

绝对收敛级数（ \(\mathbf{R}^n\) 的点的）: VII, 3, 2.

法里数列 : VII, 1, 习题 13.

几乎开集 : IX, 5, 习题 6.

波莱尔集 : IX, 6, 3.

有界集 : IX, 2, 2.

有界集（在 \(\mathbf{R}^n\) 中）: VI, I, I.

可容集 : IX, 6, 9.

贫集 : IX, 5, 2.

无处稠密集 : IX, 5, 1.

苏斯林集 : IX, 6, 2.

星形集（在 \(\mathbf{R}^n\) 中）: VI, 2, 习题 12.

瘦集 : IX, 5, 习题 2.

星形集的壳 : VI, 2, 习题 12.

立方体的面 : VI, I, I.

折线的边 : VI, 1, 习题 6.

筛 : IX, 6, 5.

筛分 : IX, 6, 5.

由筛分诱导的映射 : IX, 6, 5.

简单折线 : VI, 1, 习题 9.

角的正弦 : VIII, 2, 2.

数的正弦 : VIII, 2, 4.

截片（在线性有序集中）: IX, 附录, 习题 2.

直线的斜率 : VIII, 2, 6.

苏斯林集 : IX, 6, 2.

苏斯林空间 : IX, 6, 2.

与筛分相伴的空间 : IX, 6, 5.

贝尔空间 : IX, 5, 3.

集体正规空间 : IX, 4, 习题 18.

完全正规空间 : IX, 4, 习题 3.

完全正则空间 : IX, 1, 5.

复数空间 : VIII, 4, 1.

复射影空间 : VIII, 4, 3.

可数紧空间 : IX, 2, 习题 14.

局部仿紧空间 : IX, 4, 习题 27.

卢津空间 : IX, 6, 4.

亚紧空间 : IX, 4, 习题 25.

度量空间 : IX, 2, 1.

可数型的可度量化空间 : IX, 2, 8.

可度量化拓扑空间 : IX, 2, 5.

可度量化一致空间 : IX, 2, 4.

非贫空间 : IX, 5, 习题 7.

正规空间 : IX, 4, 1.

赋范空间 : IX, 3, 3.

闭路空间 : X, 3, 4.

道路空间 : X, 3, 4.

维数为 \(p\) 的射影线性簇所成的空间（在复射影空间中，维数为 \(n: \text{VIII}, 4, 3\) .

维数为 \(p\) 的射影线性簇所成的空间（在实射影空间中，维数为 \(n: \mathrm{VI}, 3, 5\) .

完备正规空间 : IX, 4, 习题 7.

波兰空间 : IX, 6, 1.

伪紧空间 : IX, 1, 习题 21.

实紧空间 : X, 4, 习题 17.

实数空间 : VI, I, I.

实射影空间 : VI, 3, 1.

可分可度量化空间 : IX, 2, 8.

苏斯林空间 : IX, 6, 2.

强零维空间 : IX, 6, 习题 1.

亚可度量化空间 : IX, 2, 习题 23.

完全非贫空间 : IX, 5, 习题 12.

超度量空间 : IX, 2, 习题 4.

零维空间 : IX, 6, 4.

球面 : IX, 2, 2, VI, 2, 3.

对数螺线 : VIII, 3, 习题 5.

正方形（闭的、开的） : VI, I, I.

\(\mathbf{R}^n\) 中的星形集 : VI, 2, 习题 12.

球极投影 : VI, 2, 4.

斯通–切赫紧化 : IX, 1, 习题 7.

斯通定理 : X, 4, 1.

斯通–魏尔斯特拉斯定理 : X, 4, 2.

严格筛分 : IX, 6, 5.

强零维空间 : IX, 6, 习题 1.

相伴子群 : VII, 1, 3.

亚可度量化空间 : IX, 2, 习题 23.

从属于子集族（函数族） : IX, 4, 3.

函数的支集 : IX, 4, 3.

角的正切 : VIII, 2, 2.

叉角的正切 : VIII, 2, 6.

数的正切 : VIII, 2, 4.

阿斯科利定理 : X, 2, 5.

贝尔定理 : IX, 5, 3.

迪尼定理 : X, 4, 1.

R. 埃利斯定理 : X, 3, 习题 25.

长田–斯米尔诺夫定理 : IX, 4, 习题 22.

斯通定理 : X, 4, 1.

斯通–魏尔斯特拉斯定理 : X, 4, 2.

乌雷松定理 : IX, 4, 1 与 4, 2.

瘦集 : IX, 5, 习题 2.

紧开拓扑 : X, 3, 4.

紧收敛的拓扑 : X, 1, 3 与 3, 4.

逐点收敛的拓扑 : X, 1, 3.

在 X 的子集上的逐点收敛的拓扑 : X, 1, 3.

S-收敛的拓扑 : X, 1, 2.

一致收敛的拓扑 : X, 1, 1.

在 X 的子集上的一致收敛的拓扑 : X, 1, 3.

一维环面 : V, 1, 2.

n 维环面 : VII, 1, 4.

旋转环面 : VII, 1, 习题 15.

完全非贫空间 : IX, 5, 习题 12.

正交变换 : VI, 2, 2.

三角不等式 : IX, 1, 1.

三角不等式 : VI, 2, 1.

复数的三角形式 : VIII, 2, 2.

\(m\) 个变量的三角多项式: \(\mathbf{X}\) , 4, 4.

超度量空间 : IX, 2, 习题 4.

连续实值函数用属于给定集合的连续

函数作的一致逼近 : X, 4, 1.

在 X 的子集上的一致收敛 : X, 1, 3.

在 \(\mathfrak{S}:\mathbf{X},\mathbf{1},\mathbf{2}\) 的各集合上的一致收敛

一致收敛点 : X, 1, 习题 9.

一致收敛的拓扑: X, 1, 1.

一致收敛的一致结构: X, 1, 1.

加法一致结构（ \(\mathbf{R}^n\) 的）: VI, 1, 2.

有界收敛的一致结构 : X, 1, 3.

紧收敛的一致结构 : X, 1, 3.

由伪度量族定义的一致结构 : IX, 1, 2.

由伪度量定义的一致结构 : IX, 1, 2.

有限开覆盖的一致结构 : IX, 4, 习题 17.

逐点收敛的一致结构 : X, 1, 3.

在 X 的子集上的逐点收敛的一致结构 : X, 1, 3.

预紧收敛的一致结构 : X, 1, 3.

\(\mathfrak{S}\) -收敛的一致结构 : X, 1, 2.

一致收敛的一致结构 : X, I, I.

在 X 的子集上的一致收敛的一致结构 : X, 1, 3.

在 S 的各集合上的一致收敛的一致结构 : X, 1, 2.

万有一致结构 : IX, 1, 习题 5.

一致收敛的（滤子） : X, I, I.

在 \(\mathfrak{S}\) 的每个集合上一致收敛的（映射族、滤子、序列、级数）: X, 1, 2.

一致等度连续的映射集 : X, 2, 1.

一致可和族 : X, 1, 2.

单位球、单位球面 : VI, 2, 3 与 IX, 3, 3.

角度量的单位 : VIII, 2, 3.

万有一致结构 : IX, 1, 习题 5.

乌雷松定理: IX, 4, 1 与 4, 2.

实赋值 : IX, 3, 2.

绝对值 : IX, 3, 2.

绝对值（复数的）: VIII, 1, 1.

半绝对值 : IX, 3, 习题 9.

赋值除环 : IX, 3, 2.

坐标簇 : VI, 1, 4.

方向向量（直线或射线的）: VI, 1, 4.

正交向量子空间 : VI, 2, 2.

正交向量 : VI, 2, 2.

球极投影的顶点 : VI, 2, 4.

魏尔斯特拉斯–斯通定理 : X, 4, 2.

零维空间 : IX, 6, 4
